\section{Terminal integration and the normalization constant}
\label{sec:terminal}

We now compare two squares using the same finite history.  The comparison
removes all scalar contributions from the microscopic magnetic core.
What remains is a deterministic energy difference and a small terminal
partition function.  The accuracy of the latter comparison, rather than
an estimate on the core amplitude, will give a convergent normalization.

Let $M=L^j$, $n=M/2$, and use the history with $J\leq j<2J$ as in
Proposition~\ref{prop:magnetic}.  For $Q\in\{M,LM\}$, the subscript $Q$
will indicate the free dual square of side $Q$.  Both calculations stop
at the scale in \eqref{eq:mag-common-stop}:
\begin{equation}
 v=L^h,\qquad
 j-h=\left\lceil\log_L(j^{1/10})\right\rceil,
 \qquad D=M/v.
 \label{eq:terminal-stop}
\end{equation}
In particular,
\begin{equation}
 j^{1/10}\leq D<Lj^{1/10},\qquad
 h=j-O_L(\log j).
 \label{eq:terminal-stop-size}
\end{equation}
There are $O_L(D^2)$ terminal blocks in either square.  All constants
below may depend on the parameters fixed before choosing the history,
including $L$, but not on $j$.

Write $\E_{*,Q}$ for expectation under the remaining reference Gaussian
field in the square of side $Q$, together with its common constant,
which is uniform modulo $2\pi/\alpha$.  Recall that $\alpha^2=8\pi$.
For $\varepsilon\in\{0,l\}$, define
\[
 \mathcal Z_Q^\varepsilon
 =\E_{*,Q}\mathscr Z_h(K^\varepsilon_{h,Q};\psi),
 \qquad
 \mathcal J_Q
 =\E_{*,Q}\mathscr I_h(K^l_{h,Q},I_{h,Q};\psi).
\]
Here $\mathscr Z_h$ is the compatible-support partition polynomial,
and $\mathscr I_h$ is its part with exactly one marked activity, as in
Proposition~\ref{prop:magnetic}.  In particular, $\mathcal J_Q$ does not
include the separately extracted scalar $\sigma_Q$.

\subsection{Uniform terminal integration}

The fact that $D$ grows is useful: it separates the magnetic core from
the wall at the terminal scale.  It also requires a terminal estimate
whose constants do not depend on the total number of blocks.

\begin{lemma}[Terminal partition expansion]
\label{lem:terminal-partition}
For the backgrounds of Proposition~\ref{prop:magnetic},
\begin{align}
 \mathcal Z_Q^\varepsilon&=1+O_L(D^2/j),
 \label{eq:terminal-partition-small}\\
 \log\mathcal Z_Q^\varepsilon
 &=\sum_X\E_{*,Q}K^\varepsilon_{h,Q}(X;\psi)
       +O_L(D^4/j^2).
 \label{eq:terminal-log}
\end{align}
They are real and positive for all sufficiently large $j$.  Moreover,
\begin{equation}
 |\mathcal J_Q|
 \leq C_LD^2\|I_{h,Q}\|.
 \label{eq:terminal-insertion}
\end{equation}
These estimates hold with the fixed regulator and activity-weight
reserves used in the magnetic construction.
\end{lemma}

\begin{proof}
The terminal regulator estimate in Proposition~\ref{inp:regulator} gives
\begin{equation}
 \E_{*,Q}\prod_{a=1}^p\mathcal G_h(X_a;\psi)
 \leq C^{|X_1|+\cdots+|X_p|}
 \label{eq:terminal-regulator}
\end{equation}
for compatible supports, with any one of the fixed reserved regulators
$\mathcal G_h$ used here.  The constant in this bound is independent of
$Q/v$.  To recall why the terminal estimate has this uniformity, the
positive-order differences of the remaining covariance are bounded at
scale $v$, and its energy form is bounded by the full reference form.
The blockwise Gaussian quadratic estimates in the regulator proof
therefore have uniformly bounded operator norm and trace proportional
to the number of blocks being tested.  Their exponential moments cost
at most $C^{|X|}$ on a support $X$.  Neither estimate uses an upper
bound on the number of blocks in the entire square.  Regulator
composition and monotonicity apply the same estimate to the union of
compatible supports.  The common constant causes no additional cost,
since the activities are periodic and the regulators depend only on
nonconstant fields.  These are the terminal bounds of
\cite[Sections 9 and 4]{B}; their
kernel hypotheses hold equally at the earlier stop
\eqref{eq:terminal-stop}.

Proposition~\ref{prop:magnetic} bounds each background by $C/j$ in the
activity norm.  Sum first over supports through a specified anchor
block and then over the $O_L(D^2)$ anchors.  The exponential cost in
\eqref{eq:terminal-regulator} is absorbed by the reserved activity
weight.  It follows that the absolute sum of terms of degree $p$ in
the terminal partition polynomial is bounded by
\[
                    (C_LD^2/j)^p.
\]
Since $D^2/j\to0$, the terms of degree at least two have total size
$O_L(D^4/j^2)$.  The constant term is one and the degree-one term is
the sum in \eqref{eq:terminal-log}.  This proves
\eqref{eq:terminal-partition-small}; expanding $\log(1+x)$ proves
\eqref{eq:terminal-log}.  Reality follows from the real-coefficient
construction, and closeness to one then gives positivity.  The same
calculation with one distinguished marked factor bounds the sum by
\[
 C_LD^2\|I_{h,Q}\|
       \sum_{p\geq0}(C_LD^2/j)^p.
\]
Absorbing this convergent series proves
\eqref{eq:terminal-insertion}.
\end{proof}

\subsection{Cancellation of the magnetic scalar}

Let $\widehat a_{Q/2}$ denote the ratio with the observation cutoffs
used for this history.  This notation depends on the chosen history;
only the two squares under simultaneous comparison use it below.
Set
\[
 E_Q=\|\eta_{l,Q}\|^2,
 \qquad H_Q=\exp\{-E_Q/(2\alpha^2)\}.
\]
The common preparation scalar cancels between the two sectors.  The
exact identities \eqref{eq:mag-background-products} and
\eqref{eq:mag-insertion-identity} therefore give
\begin{equation}
 \widehat a_{Q/2}
 =H_Q\rho_j
       \frac{\sigma_j\mathcal Z_Q^l+\mathcal J_Q}
            {\mathcal Z_Q^0},
 \qquad
 \rho_j=\frac{P_Q^l}{P_Q^0}>0.
 \label{eq:terminal-exact-ratio}
\end{equation}
Both $\rho_j$ and $\sigma_j$ are exactly the same for $Q=M$ and $Q=LM$.
The equality of $\sigma_j$ is algebraic; its sign has not been assumed.

\begin{lemma}[Relative removal of the residual insertion]
\label{lem:terminal-scalar}
There is $c>0$ such that, for all sufficiently large $j$,
$\sigma_j>0$ and
\begin{equation}
 \log a_{Q/2}
 =-\frac{E_Q}{2\alpha^2}+\log\rho_j+\log\sigma_j
    +\log\mathcal Z_Q^l-\log\mathcal Z_Q^0+O(e^{-cj})
 \label{eq:terminal-log-ratio}
\end{equation}
for both $Q=M$ and $Q=LM$.
\end{lemma}

\begin{proof}
The bounds on $s$, the entry index $i$, and the terminal index $h$ give
$s=O_L(j^3)$, $i=O_L(\log j)$, and $j-h=O_L(\log j)$.  By
\eqref{eq:mag-insertion-bound} and
\eqref{eq:terminal-insertion}, the marked term in the physical ratio
\eqref{eq:terminal-exact-ratio} is bounded by
\begin{equation}
 \left|H_Q\rho_j\frac{\mathcal J_Q}{\mathcal Z_Q^0}\right|
 \leq \operatorname{poly}(j)L^{-2j/3}e^{\epsilon_2j}.
 \label{eq:terminal-physical-residual}
\end{equation}
Here $H_Q\leq1$, the denominator is bounded below by
Lemma~\ref{lem:terminal-partition}, and
Proposition~\ref{prop:magnetic} gives
$|\log\rho_j|\leq\epsilon_2j$.  The positive number $\epsilon_2$ may
be chosen arbitrarily small by shrinking the background tube after
fixing the map constants.  Choose $\epsilon_2<(\log L)/6$.

Lemma~\ref{inp:regular-variation} gives
$a_{Q/2}\geq c_L L^{-j/2}$.  Consequently,
\eqref{eq:terminal-physical-residual} is exponentially small relative
to $a_{Q/2}$.  The cutoff error is even smaller: the preparation bound
has the form
\[
 C(Q/s)^2e^{-c_1s}=\exp\{-c_1s+O_L(j)\},
\]
and $s\geq c_Lj^3$.  The ordinary cutoff denominator is
$1+o(1)$ after division by its uncut partition function, so the same
bound, with a changed constant, controls the error in the ratio.
Dividing by the crude lower bound again shows that
$\widehat a_{Q/2}=a_{Q/2}(1+O(e^{-cj}))$.

The exact identity \eqref{eq:terminal-exact-ratio} now implies
\[
 H_Q\rho_j\sigma_j\frac{\mathcal Z_Q^l}{\mathcal Z_Q^0}
       =a_{Q/2}(1+O(e^{-cj})).
\]
Every factor on the left except $\sigma_j$ is positive.  Thus
$\sigma_j>0$ for large $j$, and taking logarithms proves
\eqref{eq:terminal-log-ratio}.  This argument uses neither positivity
of individual activities nor an a priori lower bound on the extracted
insertion scalar.
\end{proof}

The two scalar logarithms in \eqref{eq:terminal-log-ratio} cancel when
the square is enlarged.  It remains to compute the difference of the
terminal background terms to a summable accuracy.

\subsection{The first-order magnetic difference}

Define the first-order term
\[
 F_Q^\varepsilon
   =\sum_X\E_{*,Q}K^\varepsilon_{h,Q}(X;\psi)
\]
and its magnetic double difference
\[
 \Delta_j=(F_{LM}^l-F_{LM}^0)-(F_M^l-F_M^0).
\]
The word ``double'' here refers simply to taking first a difference
between the magnetic and ordinary sectors and then a difference
between the two square sizes.

\begin{lemma}[First-order terminal comparison]
\label{lem:terminal-first-order}
With the running gradient coordinate $t_h$ of
Proposition~\ref{prop:bulk-window},
\begin{equation}
 \Delta_j
 =\frac{t_h}{2\alpha^2}(E_{LM}-E_M)
     +O_L\left(\frac{D^2}{j^2}+\frac1{Dj}\right).
 \label{eq:terminal-first-order}
\end{equation}
\end{lemma}

\begin{proof}
We split the supports into shift-good bulk supports, shift-good wall
supports, and the remaining supports near the mark.  These are the
same support classifications used for the copying prescriptions in
Proposition~\ref{prop:magnetic}.  For a support of size $p$, all the
enlarged rectangles and test collars used below have diameter at most
$C_Lpv$.  Polynomial factors in $p$ are harmless in sums with the
reserved exponential activity weight.

First discard supports with $p>c_2D$, where $c_2>0$ is a sufficiently
small fixed constant depending on the geometry.  Their full norm is
$O(1/j)$, so their total integrated contribution in either square is
\[
                   O_L(D^2j^{-1}e^{-cD}).
\]
This is smaller than $1/(Dj)$ for all sufficiently large $j$.  We can
therefore assume $p\leq c_2D$ throughout the local comparisons.

\emph{Shift-good bulk supports.}
The pierced activity is the ordinary plane activity evaluated at
$\psi+\lambda_Q$, where the local lift satisfies
$\partial\lambda_Q=\eta_{l,Q}/\alpha$.  The remainder after the
running gradient and fundamental singletons has norm $O(j^{-2})$ by
Proposition~\ref{prop:bulk-window}.  Its translated norm has the same
bound with the fixed reserves, and summing over the $O_L(D^2)$
possible anchors costs $O_L(D^2/j^2)$.

The fundamental singleton has zero expectation, before and after the
translation, because integration over the common constant annihilates
every nonzero charge.  For the gradient singleton, the cross term has
zero Gaussian expectation.  Its translated-minus-untranslated
expectation is therefore exactly the deterministic gradient energy of
$\lambda_Q$, with coefficient $t_h/2$.  The energy allocations
$e_B^0$ in the bulk map partition this energy.  Summing the bulk
singleton contributions gives the corresponding portion of
$t_hE_Q/(2\alpha^2)$.

We may complete this portion to the full energy when taking the
difference of squares.  The omitted wall positions occupy a strip of
width $O_L(v)$ and length $O_L(Q)$.  The source bounds in
Equation~\eqref{inp:source-lifts} give $|\eta_{l,Q}|\leq C_L/M$ there,
hence an omitted energy $O_L(v/M)=O_L(D^{-1})$.  The positions omitted
near the mark can be chosen identically relative to the center in the
two squares.  Their energies individually need not be small.  Their
\emph{difference}, however, is small.  On their region of radius
$O_L(v)$, the square-to-plane alignment estimate
\eqref{eq:lift-derivatives}, together with
\eqref{inp:source-lifts}, gives
\[
 |\eta_{l,LM}-\eta_{l,M}|\leq C_L/M,
 \qquad
 |\eta_{l,LM}|+|\eta_{l,M}|
       \leq \frac{C_L}{1+\operatorname{dist}(\cdot,0)}.
\]
Thus, summing over lattice annuli,
\[
 \sum_{\operatorname{dist}(e,0)\leq C_Lv}
 \bigl||\eta_{l,LM}(e)|^2-|\eta_{l,M}(e)|^2\bigr|
 \leq C_L\left(\frac vM+\frac{v^2}{M^2}\right)
 =O_L(D^{-1}).
\]
The same estimate covers the allocation collars and half-edge
conventions.  Since $t_h=O(1/j)$, the resulting error is $O_L(1/(Dj))$.

\emph{Shift-good wall supports.}
For these supports one must use the full $O(1/j)$ activity norm; it is
not a bulk remainder.  The support is a macroscopic distance from the
source.  Subtracting a constant from the local lift, which has no
effect after common-mean integration, gives scaled test norm at most
$C_Lp/D$ on its testing neighborhood.  The source derivative bounds,
with the reflection stencils at a free wall, give this estimate also
for the derivatives required in that norm.

Let $K$ be the corresponding ordinary wall activity and let
$\lambda$ be this constant-subtracted lift.  The function
\[
                    f(u)=\E_{*,Q}K(\psi+u\lambda)
\]
is even in $u$: the ordinary activity is even in the field and the
Gaussian field, including its uniform constant, is invariant under
sign reversal.  Consequently $f'(0)=0$.  Taylor's formula, the
analytic derivative bound, and terminal regulator integration give
\[
 |f(1)-f(0)|
 \leq \frac{C_Lp^2}{D^2}\,\frac1j
       \times\text{the summable support weight}.
\]
Choose $c_2$ small enough that the translation stays inside the fixed
analytic radius; the reserved regulators handle its real part.
There are $O_L(D)$ wall anchors, with at most additional polynomial
factors in $p$ from the collars.  The weighted sum is therefore
$O_L(1/(Dj))$.  This estimate explains why a bound by the absolute
wall activity norm would lose the required accuracy.

\emph{Supports near the mark: deterministic translation.}
For a support which is not shift-good and has $p\leq c_2D$, the
enlarged support lies in the interior of both squares.  The magnetic
activities copy a common pierced-plane activity after translation by
$\gamma_Q$, and the ordinary activities copy the ordinary-plane
activity.  The change $\gamma_{LM}-\gamma_M$, after subtraction of a
constant, has local scaled test norm at most $C_Lp/D$ by
Equation~\eqref{eq:lift-derivatives}.  The first derivative estimate therefore
bounds the effect of changing this deterministic translation by
$C_Lp/(jD)$ times a summable support weight.  Such a support has only
polynomially many possible anchor positions in $p$ near the mark.
Its weighted sum is $O_L(1/(Dj))$.

\emph{Supports near the mark: terminal covariance.}
The preceding comparison must also change the terminal Gaussian law.
We give this step explicitly.  Average each activity over its common
constant first, so the resulting function is invariant under addition
of any constant to the field.  Continue the polynomial covariance
decomposition from scale $v$ to a fixed sufficiently small dyadic
fraction $c_3M$ of the smaller square.  With $c_3$ and then $c_2$
chosen using the interior clearance, all these polynomial pieces agree
on the enlarged support in the two squares, modulo common constants.
This is the exact finite-range copying property, including polynomial
continuation on the zero mode.  Common constants are invisible to the
mean-averaged activity.

We can consequently realize the restrictions of the two fields as a
common Gaussian part plus respective independent tails.  Anchor each
tail at a fixed point in the enlarged rectangle.  The positive-order
terminal covariance estimates at scale $c_3M$ imply, for every fixed
$r\geq2$,
\begin{equation}
 \bigl\|\text{anchored tail}\bigr\|_{L^r(\text{scaled test norm})}
       \leq \frac{C_{L,r}\operatorname{poly}(p)}D.
 \label{eq:terminal-tail}
\end{equation}
For completeness, the difference estimates through order five bound
the scaled differences of orders one and two in supremum on each
$v$-block by the rescaled discrete Sobolev estimate used in the
regulator proof.  Summing the first differences along paths from the
anchor bounds the field value itself.  The enclosing rectangle has
block dimensions $O_L(1+p)$, so these sums introduce only a polynomial
factor in $p$.  Gaussian moment estimates give the stated $L^r$
bounds.  The factor $v/M=D^{-1}$ comes from the first positive-order
tail estimate.

Interpolate each tail separately to zero, keeping the common part
fixed.  The fundamental theorem of calculus and the first derivative
activity bound reduce the difference of expectations to
\eqref{eq:terminal-tail}.  H\"older's inequality with a fixed terminal
regulator reserve bounds the accompanying regulator.  Along this
interpolation the covariance on gradients is no larger than the
original one, so the terminal integration estimate remains valid.
The deterministic alignment translations cost only another
exponential-in-$p$ factor.  The local variational regulator agrees in
the two squares by finite-range copying on the open testing support;
all remaining exponential factors are absorbed by the reserved
weights.  Summing the $O(1/j)$ full activity norms, polynomially many
anchors, and exponentially weighted shapes gives $O_L(1/(Dj))$.

This compares covariances only through the small smooth tails that the
activity can detect.  No comparison of microscopic-dimensional
Gaussian densities is needed.  Combining the bulk, wall, and
near-mark estimates proves \eqref{eq:terminal-first-order}.
\end{proof}

\subsection{The deterministic energy increment}

We next compute the constant in the leading energy increment.  Its
normalization is important: the Laplacian here counts each undirected
unit edge once.

\begin{lemma}[Square Green-function increment]
\label{lem:terminal-green}
Let $G_n$ be the Green function of the unscaled Laplacian killed on
the boundary vertices of the spin square $\Lambda_n$.  For the
connections in the dual squares of sides $M=2n$ and $LM$, there is
$c_4>0$ such that
\begin{equation}
 E_{LM}-E_M
  =2\pi\log L+O_L(n^{-c_4}).
 \label{eq:terminal-energy}
\end{equation}
\end{lemma}

\begin{proof}
The orthogonal energy minimization and crossing correspondence in
Equation~\eqref{inp:connection-energy} give the exact identity
\begin{equation}
                E_M=(2\pi)^2G_n(0,0).
 \label{eq:terminal-green-energy}
\end{equation}
We verify the accuracy needed for the Green increment:
\begin{equation}
 G_{Ln}(0,0)-G_n(0,0)
       =\frac{\log L}{2\pi}+O_L(n^{-c_4}).
 \label{eq:terminal-green-increment}
\end{equation}

Write $M=2n$, fix a sufficiently small dyadic stop fraction $\rho$,
and use the low-pass multiplier
\[
                    q_t(\lambda)=P_t(1-\lambda/64)
\]
from Equation~\eqref{inp:polynomials}, now applied to the killed
Laplacian.  Odd reflection gives the same binary decomposition and
multiplier estimates in this geometry, as in
\cite[Lemmas 4.1 and 4.2]{B}.
Let $d_t$ be the plane diagonal at the origin of the polynomial
covariance $A^{-1}(q_t(A)-q_{2t}(A))$, with polynomial continuation
at zero, and let $T_Q$ be the smoothed terminal diagonal at the center
of the square of side $Q$, stopped at scale $\rho Q$.
All preterminal pieces at the center are plane pieces by their finite
range.  Their common terms cancel exactly, giving
\begin{equation}
 G_{Ln}(0,0)-G_n(0,0)
   =\sum_{r=0}^{\log_2L-1}d_{2^r\rho M}+T_{LM}-T_M.
 \label{eq:terminal-green-telescope}
\end{equation}
The plane diagonal estimate for each of these $\log_2L$ terms is
$\log2/(2\pi)$ with an error bounded by a negative power of its
scale.  This is the continuum kernel estimate in
\cite[Lemma 8.1]{B}; the value of the limiting diagonal is
the radial integral
\[
 \frac1{2\pi}\int_0^\infty
      \frac{P(r)-P(2r)}r\,dr=\frac{\log2}{2\pi},
\]
where $P$ is the continuum low-pass multiplier, $P(0)=1$, and
$P(r)\to0$ at infinity.  Since their scales are comparable to
$n$, their total contribution is
$\log L/(2\pi)+O_L(n^{-c})$.

The two smoothed terminal diagonals have the same continuum limit,
with power accuracy.  We verify this stronger quantitative assertion
directly.  The eigenvalues of the killed square are
\[
 \lambda_{k,M}
   =4\sin^2\frac{\pi k_1}{2M}
       +4\sin^2\frac{\pi k_2}{2M},
 \qquad 1\leq k_1,k_2\leq M-1.
\]
The smoothed terminal diagonal at the center is consequently
\begin{equation}
 T_M=\frac4{M^2}
     \sum_{k_1,k_2=1}^{M-1}
     \sin^2\frac{\pi k_1}{2}\sin^2\frac{\pi k_2}{2}
     \frac{q_{\rho M}(\lambda_{k,M})}{\lambda_{k,M}}.
 \label{eq:terminal-sine-sum}
\end{equation}
The elementary eigenvalue bounds and the multiplier estimate give
\[
 M^2\lambda_{k,M}\geq c|k|^2,
 \qquad
 |q_{\rho M}(\lambda_{k,M})|
       \leq C_\rho\min\{1,|k|^{-2p}\},
\]
where $p$ is the fixed power in the kernel construction.  Thus the
part of \eqref{eq:terminal-sine-sum} with $|k|>B$ is bounded by
$C_\rho B^{-c}$ for some $c>0$, uniformly in $M$.

For $|k|\leq B$, one has
$M^2\lambda_{k,M}=\pi^2|k|^2+O(B^4/M^2)$.  In the sine formula for
$q_{\rho M}$, the rescaled numerator argument and denominator converge
to their continuum values with errors bounded by a polynomial in $B$
times $M^{-1}$.  The ratio is interpreted continuously at zero.
Its fixed power, the reciprocal eigenvalue, and the $O(B^2)$ terms
in the sum therefore have combined error at most
$C_\rho B^{C'}/M$ for a fixed $C'>0$.  This is a deliberately weaker
bound than the eigenvalue Taylor error, and suffices here.  The
limiting sum depends on $\rho$ but not on $M$.  Taking $B=n^\delta$
with a sufficiently small fixed $\delta>0$ makes both errors negative
powers of $n$.  The stop fractions are the same in the two squares,
so their limiting terminal sums cancel.  This proves
\eqref{eq:terminal-green-increment}.  Multiplying it by $(2\pi)^2$
in \eqref{eq:terminal-green-energy} proves
\eqref{eq:terminal-energy}.
\end{proof}

\subsection{A summable ratio and the full limit}

All quantities involving the microscopic magnetic core have now
cancelled.  The universal logarithmic correction comes from the
running gradient coefficient in Lemma~\ref{lem:terminal-first-order}.

\begin{proposition}[Summable enlargement ratio]
\label{prop:ratio}
For $n_j=L^j/2$, there exists $\epsilon>0$ such that
\begin{equation}
 \log a_{n_{j+1}}-\log a_{n_j}
   =-\frac{\log L}{8}+\frac1{16j}
       +O_L(j^{-1-\epsilon}).
 \label{eq:terminal-ratio}
\end{equation}
\end{proposition}

\begin{proof}
Subtract the two identities in
Lemma~\ref{lem:terminal-scalar}.  The common $\log\rho_j$ and
$\log\sigma_j$ cancel exactly.  Lemmas~\ref{lem:terminal-partition}
and~\ref{lem:terminal-first-order} give
\[
 \log a_{Ln}-\log a_n
 =-\frac{1-t_h}{2\alpha^2}(E_{LM}-E_M)
   +O_L\left(\frac{D^4}{j^2}+\frac1{Dj}+e^{-cj}\right).
\]
Use $\alpha^2=8\pi$ and Lemma~\ref{lem:terminal-green} to obtain
\begin{equation}
 \log a_{Ln}-\log a_n
 =-\frac{\log L}{8}
      +t_h\frac{\log L}{8}
      +O_L\left(\frac{D^4}{j^2}+\frac1{Dj}+n^{-c_5}\right)
 \label{eq:terminal-ratio-before-flow}
\end{equation}
for some $c_5>0$.  Proposition~\ref{prop:bulk-window} and
\eqref{eq:terminal-stop-size} imply
\[
 t_h=\frac1{2h\log L}+O_L\left(\frac{\log h}{h^2}\right)
     =\frac1{2j\log L}+O_L\left(\frac{\log j}{j^2}\right).
\]
Finally $D^4/j^2=O_L(j^{-8/5})$ and
$1/(Dj)=O(j^{-11/10})$.  Substitution in
\eqref{eq:terminal-ratio-before-flow} proves
\eqref{eq:terminal-ratio}, for example with any fixed
$0<\epsilon\leq1/10$.
\end{proof}

\begin{proof}[Proof of Theorem~\ref{thm:main}]
For $n>1$, define
\[
 B(n)=\log a_n+\frac18\log n-\frac1{16}\log\log n.
\]
Because $\log n_j=j\log L-\log2$,
\[
 \log\log n_{j+1}-\log\log n_j
       =\frac1j+O_L(j^{-2}).
\]
Proposition~\ref{prop:ratio} therefore implies
$B(n_{j+1})-B(n_j)=O_L(j^{-1-\epsilon})$.
The increments are absolutely summable.  Each $B(n_j)$ is a finite
real number by positivity of $a_{n_j}$, so there is a finite real
limit $B_\infty$ along this geometric sequence.

It remains to remove the restriction on the side lengths.  For an
arbitrary integer $n$ tending to infinity, let $n_j$ be the preceding
geometric size, so $n_j\leq n<n_{j+1}$.  From any sequence of such
$n$, extract a subsequence on which $n/n_j\to d\in[1,L]$.
Lemma~\ref{inp:regular-variation} gives
\[
 \log a_n-\log a_{n_j}\longrightarrow-\frac18\log d.
\]
On this subsequence, $\log n-\log n_j\to\log d$ and
$\log\log n-\log\log n_j\to0$.  Hence
$B(n)-B(n_j)\to0$.  Every sequence admits such a subsequence, and the
same argument applies to every subsequence; thus $B(n)\to B_\infty$
along all integers.  Setting $A_{\rm XY}=e^{B_\infty}$ proves the
claimed asymptotic.  The finiteness of $B_\infty$ gives both
$A_{\rm XY}>0$ and $A_{\rm XY}<\infty$.
\end{proof}
